\definecolor{dsdlcodebg}{RGB}{246,248,250}
\definecolor{dsdlcodeborder}{RGB}{208,215,222}
\newsavebox{\dsdlcodebox}
\newenvironment{dsdlcodeblock}
{%
  \par\noindent
  \begin{latin}
  \begin{lrbox}{\dsdlcodebox}
  \begin{minipage}{0.94\linewidth}
  \small
}
{%
  \end{minipage}
  \end{lrbox}
  \setlength{\fboxsep}{6pt}
  \fcolorbox{dsdlcodeborder}{dsdlcodebg}{\usebox{\dsdlcodebox}}
  \end{latin}
  \par
}

\section{مقدمه}

\لر{UART} روشی رایج برای انتقال سریال و ناهمزمان داده است. چون
ساعت\LTRfootnote{Clock} مستقلی همراه داده ارسال نمی‌شود،
فرستنده و گیرنده باید بر سر
نرخ بیت\LTRfootnote{Bit Rate} و قالب قاب\LTRfootnote{Frame Format} توافق داشته باشند.
هدف آزمایش حاضر، طراحی یک فرستنده و گیرندهٔ هفت‌بیتی قابل
سنتز، اتصال آن‌ها در مسیر بازگشت
داخلی\LTRfootnote{Internal Loopback} و اثبات عملکردشان با
میزآزمون‌های خودبررسی\LTRfootnote{Self-checking Testbenches} است.

اهداف جزئی آزمایش عبارت‌اند از تبدیل دادهٔ موازی هفت‌بیتی به
دنبالهٔ سریال\LTRfootnote{Serial Stream}، بازسازی همان
داده در گیرنده، تولید و کنترل بیت توازن، کنترل بیت
توقف، اعلام وضعیت اشتغال فرستنده، اعلام
پایان دریافت و اعتبار قاب، و تحلیل زمانی
همهٔ این رخدادها در شکل‌موج شبیه‌سازی.


\section{قالب قاب و محاسبهٔ توازن}

خط سریال در حالت بیکار\LTRfootnote{Idle State} مقدار منطقی ۱ دارد. افت خط به صفر آغاز
قاب را مشخص می‌کند. پس از بیت شروع\LTRfootnote{Start Bit}، بیت توازن و سپس بیت‌های
\کد{D0} تا \کد{D6} ارسال می‌شوند و در پایان، بیت توقف
با مقدار ۱ قرار می‌گیرد. بنابراین، حالت \لر{Idle} پیش‌شرط آغاز دریافت است و خود قاب شامل
ده بیت به ترتیب زیر است:

\begin{latin}
\begin{center}
\texttt{Idle(1) -> Start(0) -> Parity -> D0 -> D1 -> D2 -> D3 -> D4 -> D5 -> D6 -> Stop(1)}
\end{center}
\end{latin}

دادهٔ هفت‌بیتی به روش \لر{LSB-first} منتقل می‌شود؛ یعنی \کد{D0}، کم‌ارزش‌ترین
بیت، زودتر از \کد{D6}، پرارزش‌ترین
بیت، روی خط قرار می‌گیرد. بیت توازن با \لر{XOR}
کاهشی همهٔ بیت‌های داده محاسبه می‌شود:

\begin{equation}
\mathrm{Parity}=D_0 \oplus D_1 \oplus D_2 \oplus D_3 \oplus D_4 \oplus D_5 \oplus D_6
\end{equation}

برای دادهٔ نمونهٔ \کد{7'b1011001} یا \کد{0x59}، بیت‌ها در ترتیب ارسال برابر ۱، ۰، ۰،
۱، ۱، ۰، ۱ هستند. تعداد بیت‌های یک برابر چهار است؛ در نتیجه \لر{XOR} آن‌ها صفر و بیت
توازن برابر ۰ است. قاب کامل این داده پس از خط بیکار چنین توالی‌ای دارد:

\begin{latin}
\begin{center}
\texttt{Start=0, Parity=0, D0..D6=1,0,0,1,1,0,1, Stop=1}
\end{center}
\end{latin}


\section{زمان‌بندی و انتخاب پارامتر بیت}

در نرخ ۱۱۵۲۰۰ بیت بر ثانیه، زمان هر بیت برابر معکوس این نرخ، یعنی تقریباً ۸٫۶۸
میکروثانیه است. برای ساعت ۵۰ مگاهرتز، تعداد چرخهٔ ساعت در هر بیت از رابطهٔ زیر به‌دست
می‌آید:

\begin{equation}
\frac{50\,000\,000}{115\,200}=434.027\ldots \simeq 434
\end{equation}

ازاین‌رو مقدار صحیح ۴۳۴ برای \کد{BIT\_TICKS} انتخاب شده است. هر حالت زمانی فرستنده، یک
نماد را به مدت ۴۳۴ لبهٔ ساعت نگه می‌دارد و گیرنده نیز با همین مقیاس نمونه‌برداری می‌کند.
میزآزمون اختصاصی قالب قاب از همین مقدار استفاده می‌کند تا الزام واقعی ۱۱۵۲۰۰ بیت بر
ثانیه را اثبات کند.

برای شکل‌موج‌های آموزشی، \کد{BIT\_TICKS=10} انتخاب شده است. با ساعت ۱۰ نانوثانیه‌ای،
هر نماد ۱۰۰ نانوثانیه طول دارد و مرز حالت‌ها، شمارنده و نمونه‌برداری میانی به‌وضوح قابل
مشاهده‌اند. برای آزمون تجمیعی و پیمایش تمام مقادیر، مقدار ۲ زمان اجرا و طول پروندهٔ موج
را کاهش می‌دهد، بی‌آنکه ترتیب منطقی قاب تغییر کند. این سه مقدار مکمل یکدیگرند: ۴۳۴ برای
دقت نرخ، ۱۰ برای خوانایی موج و ۲ برای پوشش سریع و گسترده.


\section{معماری طرح و اتصال بازگشتی}

پروندهٔ \کد{UARTSender.v} فرستنده و ماشین
حالت ارسال را دربر دارد. پروندهٔ \کد{UARTReceiver.v}
گیرنده، همگام‌سازی ورودی،
نمونه‌برداری، کنترل توازن و کنترل بیت توقف را پیاده‌سازی می‌کند.
پروندهٔ \کد{UARTTop.v} نمونه‌ای از هر دو پیمانه\LTRfootnote{Module} می‌سازد و
\کد{tx} فرستنده را مستقیماً به \کد{rx} گیرنده متصل می‌کند. پروندهٔ \کد{Tester.v}
میزآزمون تجمیعی است و ساعت، بازنشانی\LTRfootnote{Reset}،
تحریک‌ها\LTRfootnote{Stimuli}، کنترل خودکار و پیام‌های نتیجه را تولید
می‌کند. میزآزمون‌های مستقل نیز هر رفتار مهم را جداگانه و با شکل‌موجی خوانا بررسی می‌کنند.

اتصال سطح بالا در خطوط ۱۸ تا ۳۴ پروندهٔ \کد{UARTTop.v} آمده است. همان خروجی
\کد{tx} هم به بیرون پیمانه ارائه می‌شود و هم ورودی \کد{rx} نمونهٔ گیرنده است:

\begin{dsdlcodeblock}
\begin{verbatim}
18  UARTSender #(.BIT_TICKS(BIT_TICKS)) sender (
19      .clk(clk),
20      .rstN(rstN),
21      .new_data(new_data),
22      .send_data(send_data),
23      .tx(tx),
24      .busy(busy)
25  );
26
27  UARTReceiver #(.BIT_TICKS(BIT_TICKS)) receiver (
28      .clk(clk),
29      .rstN(rstN),
30      .rx(tx),
31      .rec_data(rec_data),
32      .rec_new_data(rec_new_data),
33      .correct_data(correct_data)
34  );
\end{verbatim}
\end{dsdlcodeblock}

بنابراین، هر دادهٔ پذیرفته‌شده در فرستنده باید پس از گذشت یک قاب کامل در خروجی گیرنده
دیده شود. این بازگشت داخلی، صحت هم‌زمان دو ماشین حالت را بدون نیاز به یک منبع سریال
خارجی آزمایش می‌کند.


\section{طراحی فرستنده}

ورودی‌های فرستنده \کد{clk}، \کد{rstN}، \کد{new\_data} و
\کد{send\_data[6:0]} و خروجی‌های آن \کد{tx} و \کد{busy} هستند.
بازنشانی فعال-پایین\LTRfootnote{Active-Low}، حالت را به \لر{IDLE}،
شمارنده‌ها و ثبات‌ها را به صفر و خط
\کد{tx} را به
حالت بیکار منطقی ۱ می‌برد.

ماشین حالت فرستنده پنج حالت \لر{IDLE}، \لر{START}، \لر{PARITY}، \لر{DATA} و
\لر{STOP} دارد. خطوط ۴۵ تا ۴۸ ارتباط حالت با دو خروجی اصلی را نشان می‌دهند:

\begin{dsdlcodeblock}
\begin{verbatim}
45  assign busy = state != IDLE;
46  assign tx = state == START  ? 1'b0 :
47              state == PARITY ? parity_latch :
48              state == DATA   ? data_latch[bit_index] : 1'b1;
\end{verbatim}
\end{dsdlcodeblock}

سیگنال \کد{busy} از نامساوی‌بودن حالت جاری با \لر{IDLE} ساخته می‌شود؛ بنابراین، از ورود
به \لر{START} تا پایان \لر{STOP} فعال می‌ماند. در \لر{START} مقدار صفر، در
\لر{PARITY} مقدار \کد{parity\_latch} و در \لر{DATA} بیت انتخاب‌شده با
\کد{bit\_index} روی \کد{tx} قرار می‌گیرد. در \لر{STOP} و \لر{IDLE} خروجی ۱ است.

خطوط ۵۹ تا ۱۰۸، قفل‌کردن ورودی و حرکت اصلی ماشین حالت را نشان می‌دهند:

\begin{dsdlcodeblock}
\begin{verbatim}
59  IDLE: begin
60      tick_count <= {TICK_WIDTH{1'b0}};
61      bit_index  <= 3'd0;
62      if (new_data) begin
63          data_latch   <= send_data;
64          parity_latch <= ^send_data;
65          state        <= START;
66      end
67  end
68
69  START: begin
70      if (tick_count == LAST_TICK) begin
71          tick_count <= {TICK_WIDTH{1'b0}};
72          state <= PARITY;
73      end else begin
74          tick_count <= tick_count + 1'b1;
75      end
76  end
77
\end{verbatim}
\end{dsdlcodeblock}

ادامهٔ ماشین، بیت توازن و سپس هفت بیت داده را مدیریت می‌کند:

\begin{dsdlcodeblock}
\begin{verbatim}
78  PARITY: begin
79      if (tick_count == LAST_TICK) begin
80          tick_count <= {TICK_WIDTH{1'b0}};
81          bit_index <= 3'd0;
82          state <= DATA;
83      end else begin
84          tick_count <= tick_count + 1'b1;
85      end
86  end
87
88  DATA: begin
89      if (tick_count == LAST_TICK) begin
90          tick_count <= {TICK_WIDTH{1'b0}};
91          if (bit_index == 3'd6) begin
92              state <= STOP;
93          end else begin
94              bit_index <= bit_index + 1'b1;
95          end
96      end else begin
97          tick_count <= tick_count + 1'b1;
98      end
99  end
100
\end{verbatim}
\end{dsdlcodeblock}

قطعهٔ پایانی، بیت توقف را نگه می‌دارد و فرستنده را به حالت بیکار بازمی‌گرداند:

\begin{dsdlcodeblock}
\begin{verbatim}
101  STOP: begin
102      if (tick_count == LAST_TICK) begin
103          tick_count <= {TICK_WIDTH{1'b0}};
104          state <= IDLE;
105      end else begin
106          tick_count <= tick_count + 1'b1;
107      end
108  end
\end{verbatim}
\end{dsdlcodeblock}

در \لر{IDLE}، تنها هنگام فعال‌بودن \کد{new\_data}، ورودی در \کد{data\_latch} و
توازن آن در \کد{parity\_latch} ذخیره می‌شود. این قفل‌کردن\LTRfootnote{Latching} باعث
می‌شود تغییر بعدی
\کد{send\_data} قاب در حال ارسال را عوض نکند. منطق پذیرش نیز فقط در شاخهٔ \لر{IDLE}
قرار دارد؛ در نتیجه، درخواست جدید در زمان \کد{busy} بدون اثر است. شمارندهٔ
\کد{tick\_count} مدت هر نماد را اندازه می‌گیرد و \کد{bit\_index} در حالت داده از صفر تا
شش افزایش می‌یابد تا ترتیب \لر{LSB-first} مستقیماً از
\کد{data\_latch[bit\_index]} حاصل شود.


\section{طراحی گیرنده}

ورودی‌های گیرنده \کد{clk}، \کد{rstN} و \کد{rx} و خروجی‌های آن
\کد{rec\_data[6:0]}، \کد{rec\_new\_data} و \کد{correct\_data} هستند. چون
\کد{rx} نسبت به \کد{clk} ناهمزمان است، ابتدا از دو
فلیپ‌فلاپ پیاپی عبور می‌کند. خطوط ۴۹ تا ۵۸ این
همگام‌ساز را نشان می‌دهند:

\begin{dsdlcodeblock}
\begin{verbatim}
49  // Two flip-flops isolate the FSM from an asynchronous serial input.
50  always @(posedge clk or negedge rstN) begin
51      if (!rstN) begin
52          rx_meta <= 1'b1;
53          rx_sync <= 1'b1;
54      end else begin
55          rx_meta <= rx;
56          rx_sync <= rx_meta;
57      end
58  end
\end{verbatim}
\end{dsdlcodeblock}

\کد{rx\_meta} نخستین نمونه و \کد{rx\_sync} نمونهٔ همگام‌شده‌ای است که ماشین حالت
مصرف می‌کند. این ساختار احتمال رسیدن متاستابیلیتی\LTRfootnote{Metastability} به منطق
کنترلی را کاهش می‌دهد و تأخیر
دو لبهٔ ساعتی قابل پیش‌بینی ایجاد می‌کند.

گیرنده نیز حالت‌های \لر{IDLE}، \لر{START}، \لر{PARITY}، \لر{DATA} و
\لر{STOP} دارد. خطوط ۷۱ تا ۱۲۹ اعلان‌های پالسی، تأیید میانی شروع، نمونه‌برداری و کنترل
نهایی قاب را پیاده‌سازی می‌کنند:

\begin{dsdlcodeblock}
\begin{verbatim}
71  // Completion and validity are aligned one-clock notifications.
72  rec_new_data <= 1'b0;
73  correct_data <= 1'b0;
74
75  case (state)
76      IDLE: begin
77          tick_count <= {TICK_WIDTH{1'b0}};
78          bit_index <= 3'd0;
79          if (!rx_sync)
80              state <= START;
81      end
82
83      START: begin
84          if (tick_count == HALF_TICK) begin
85              tick_count <= {TICK_WIDTH{1'b0}};
86              state <= rx_sync ? IDLE : PARITY;
87          end else begin
88              tick_count <= tick_count + 1'b1;
89          end
90      end
91
\end{verbatim}
\end{dsdlcodeblock}

در ادامه، بیت توازن ذخیره و هفت بیت داده به‌ترتیب نمونه‌برداری می‌شوند:

\begin{dsdlcodeblock}
\begin{verbatim}
92  PARITY: begin
93      if (tick_count == LAST_TICK) begin
94          tick_count <= {TICK_WIDTH{1'b0}};
95          bit_index <= 3'd0;
96          parity_latch <= rx_sync;
97          state <= DATA;
98      end else begin
99          tick_count <= tick_count + 1'b1;
100      end
101  end
102
103  DATA: begin
104      if (tick_count == LAST_TICK) begin
105          tick_count <= {TICK_WIDTH{1'b0}};
106          data_latch[bit_index] <= rx_sync;
107          if (bit_index == 3'd6) begin
108              state <= STOP;
109          end else begin
110              bit_index <= bit_index + 1'b1;
111          end
112      end else begin
113          tick_count <= tick_count + 1'b1;
114      end
115  end
116
\end{verbatim}
\end{dsdlcodeblock}

قطعهٔ پایانی، بیت توقف و برابری توازن را بررسی و خروجی‌های پالسی را تولید می‌کند:

\begin{dsdlcodeblock}
\begin{verbatim}
117  STOP: begin
118      if (tick_count == LAST_TICK) begin
119          tick_count <= {TICK_WIDTH{1'b0}};
120          rec_new_data <= 1'b1;
121          if (rx_sync && ((^data_latch) == parity_latch)) begin
122              rec_data <= data_latch;
123              correct_data <= 1'b1;
124          end
125          state <= IDLE;
126      end else begin
127          tick_count <= tick_count + 1'b1;
128      end
129  end
\end{verbatim}
\end{dsdlcodeblock}

در \لر{IDLE}، صفرشدن \کد{rx\_sync} شروع احتمالی قاب است. ماشین در \لر{START} تا
نیمهٔ بیت صبر می‌کند؛ اگر خط در \کد{HALF\_TICK} دوباره ۱ شده باشد، رخداد یک شروع کاذب
تلقی می‌شود و ماشین به \لر{IDLE} بازمی‌گردد. در غیر این صورت، بیت توازن در
\کد{parity\_latch} و هفت نمونهٔ بعدی از \کد{D0} تا \کد{D6} در
\کد{data\_latch} ذخیره می‌شوند.

در پایان \لر{STOP}، \کد{rec\_new\_data} برای دقیقاً یک چرخهٔ ساعت فعال می‌شود و پایان
پردازش قاب را اعلام می‌کند. شرط اعتبار هم‌زمان شامل \کد{rx\_sync=1} برای بیت توقف و
برابری \کد{parity\_latch} با \لر{XOR} کاهشی \کد{data\_latch} است. فقط در صورت
برقرار بودن هر دو شرط، \کد{rec\_data} با دادهٔ تازه به‌روز و \کد{correct\_data} نیز
برای همان یک چرخه فعال می‌شود. در قاب نامعتبر، اعلان پایان همچنان ایجاد می‌شود، ولی
\کد{correct\_data} صفر می‌ماند و \کد{rec\_data} آخرین دادهٔ معتبر را حفظ می‌کند.


\section{روش میزآزمون و پوشش آزمون‌ها}

ساعت اصلی در میزآزمون‌ها دورهٔ ۱۰ نانوثانیه دارد. در آغاز، \کد{rstN} صفر می‌شود و سپس
مدار در یک لبهٔ کنترل‌شده از بازنشانی خارج می‌گردد. درخواست ارسال با قرار دادن داده روی
\کد{send\_data} و فعال‌کردن \کد{new\_data} برای یک چرخهٔ ساعت ساخته می‌شود.
کنترل‌گرهای خودکار با انتظار برای \کد{busy} و \کد{rec\_new\_data}، مقدار
\کد{rec\_data}، اعتبار \کد{correct\_data} و پهنای
پالس\LTRfootnote{Pulse Width} را بررسی می‌کنند. برای قاب‌های خراب، یک محرک
خام\LTRfootnote{Raw Driver} \کد{rx} بیت‌های قاب را مستقیماً تولید می‌کند تا توازن یا بیت
توقف عمداً تغییر داده شود.

سیگنال‌های اصلی قرارگرفته در پنجرهٔ \لر{Wave} عبارت‌اند از \کد{clk}، \کد{rstN}،
\کد{new\_data}، \کد{send\_data}، \کد{tx}، \کد{busy}، \کد{rec\_data}،
\کد{rec\_new\_data} و \کد{correct\_data}. در آزمون‌های اختصاصی، سیگنال‌های داخلی
مانند \کد{state}، \کد{tick\_count}، \کد{bit\_index}، \کد{data\_latch}،
\کد{parity\_latch}، \کد{rx\_meta} و \کد{rx\_sync} نیز برای تحلیل دقیق‌تر افزوده
شده‌اند.

جدول~\رجوع{جدول:پوشش آزمون‌ها} محرک و نتیجهٔ مورد انتظار همهٔ سناریوهای الزامی و تکمیلی را
خلاصه می‌کند.

\شروع{لوح}[H]
\تنظیم‌ازوسط
\شرح{ماتریس پوشش میزآزمون‌های \لر{UART}}

\شروع{جدول}{|c|p{4cm}|p{8.2cm}|}
\خط‌پر
\سیاه ردیف & \سیاه سناریو و محرک & \سیاه نتیجهٔ مورد انتظار \\
\خط‌پر \خط‌پر
۱ & دریافت صحیح \کد{0x59} & \کد{rec\_data=0x59} و پالس هم‌زمان پایان و اعتبار \\
\خط‌پر
۲ & قالب دقیق با \کد{BIT\_TICKS=434} & ده بیت قاب، هرکدام به مدت ۴۳۴ چرخه \\
\خط‌پر
۳ & \کد{0x2D} با توازن معکوس ۱ & پایان قاب، اعتبار صفر و حفظ دادهٔ معتبر قبلی \\
\خط‌پر
۴ & \کد{0x2D} با \کد{Stop=0} & رد قاب حتی با توازن صحیح \\
\خط‌پر
۵ & \کد{0x59} و سپس \کد{0x2D} & دو دریافت معتبر و مرتب در نخستین فرصت قانونی \\
\خط‌پر
۶ & درخواست \کد{0x2D} هنگام ارسال \کد{0x59} & قاب جاری بدون تغییر و نبود دریافت دوم \\
\خط‌پر
۷ & بازنشانی با \کد{send\_data=0x7F} & خط بیکار، \کد{busy=0} و خروجی‌های پاک \\
\خط‌پر
۸ & قاب معتبر و قاب با توازن خراب & پالس یک‌چرخه‌ای پایان برای هر دو و اعتبار فقط برای قاب سالم \\
\خط‌پر
۹ & پیمایش \کد{0x00} تا \کد{0x7F} & تطابق و اعتبار برای هر ۱۲۸ مقدار \\
\خط‌پر
۱۰ & پالس شروع ۱۰ نانوثانیه‌ای و سپس \کد{0x59} & رد شروع کاذب و بازیابی کامل \\
\خط‌پر
\پایان{جدول}

\برچسب{جدول:پوشش آزمون‌ها}
\پایان{لوح}

میزآزمون \کد{tb\_uart\_frame\_format.v} مقدار \کد{BIT\_TICKS} را در خط ۵ برابر ۴۳۴
و دادهٔ آزمون را در خط ۶ برابر \کد{7'b1011001} قرار می‌دهد. تابع کنترل خطوط ۲۸ تا ۴۰،
هر نماد را در تمام ۴۳۴ چرخه می‌سنجد:

\begin{dsdlcodeblock}
\begin{verbatim}
5  localparam integer BIT_TICKS = 434;
6  localparam [6:0] TEST_DATA = 7'b1011001;
28  task check_symbol;
29      input expected_tx;
30      integer cycle;
31      begin
32          for (cycle = 0; cycle < BIT_TICKS; cycle = cycle + 1) begin
33              if (tx !== expected_tx || busy !== 1'b1)
34                  $fatal(1,
35                         "FAIL: symbol=%b cycle=%0d tx=%b busy=%b",
36                         expected_tx, cycle, tx, busy);
37              @(negedge clk);
38          end
39      end
40  endtask
\end{verbatim}
\end{dsdlcodeblock}

پس از اعمال بازنشانی و پالس یک‌چرخه‌ای \کد{new\_data}، خطوط ۵۶ تا ۷۱ بیت شروع، توازن،
هفت بیت داده و بیت توقف را به همان ترتیب کنترل می‌کنند و بازگشت نهایی به حالت بیکار را
می‌سنجند:

\begin{dsdlcodeblock}
\begin{verbatim}
56  @(negedge clk);
57  send_data = TEST_DATA;
58  new_data = 1'b1;
59  @(negedge clk);
60  new_data = 1'b0;
61
62  check_symbol(1'b0);
63  check_symbol(^TEST_DATA);
64  for (bit_index = 0; bit_index < 7; bit_index = bit_index + 1)
65      check_symbol(TEST_DATA[bit_index]);
66  check_symbol(1'b1);
67
68  if (tx !== 1'b1 || busy !== 1'b0)
69      $fatal(1, "FAIL: sender did not return to idle");
70
71  $display("PASS: sender frame format and 434-cycle timing");
\end{verbatim}
\end{dsdlcodeblock}

پایش پالس‌های خروجی در خطوط ۶۲ تا ۹۳ پروندهٔ \کد{Tester.v} انجام می‌شود. علاوه بر ثبت
هر تکمیل، این بخش خطا می‌دهد اگر \کد{rec\_new\_data} بیش از یک چرخه بماند یا
\کد{correct\_data} خارج از چرخهٔ تکمیل فعال شود:

\begin{dsdlcodeblock}
\begin{verbatim}
62  // Record every completion and enforce the one-cycle output contract.
63  always @(negedge clk or negedge rstN) begin
64      if (!rstN) begin
65          top_completion_count = 0;
66          raw_completion_count = 0;
67          previous_top_new = 1'b0;
68          previous_raw_new = 1'b0;
69      end else begin
70          if (rec_new_data && previous_top_new)
71              $fatal(1, "FAIL: loopback rec_new_data exceeded one cycle");
72          if (raw_rec_new_data && previous_raw_new)
73              $fatal(1, "FAIL: raw rec_new_data exceeded one cycle");
74          if (correct_data && !rec_new_data)
75              $fatal(1, "FAIL: loopback correct_data is not completion-aligned");
76          if (raw_correct_data && !raw_rec_new_data)
77              $fatal(1, "FAIL: raw correct_data is not completion-aligned");
78
\end{verbatim}
\end{dsdlcodeblock}

بخش دوم پایش، داده و اعتبار هر تکمیل را ثبت و وضعیت چرخهٔ پیشین را نگه‌داری می‌کند:

\begin{dsdlcodeblock}
\begin{verbatim}
79          if (rec_new_data) begin
80              top_data_history[top_completion_count] = rec_data;
81              top_correct_history[top_completion_count] = correct_data;
82              top_completion_count = top_completion_count + 1;
83          end
84          if (raw_rec_new_data) begin
85              raw_data_history[raw_completion_count] = raw_rec_data;
86              raw_correct_history[raw_completion_count] = raw_correct_data;
87              raw_completion_count = raw_completion_count + 1;
88          end
89
90          previous_top_new = rec_new_data;
91          previous_raw_new = raw_rec_new_data;
92      end
93  end
\end{verbatim}
\end{dsdlcodeblock}

سناریوهای خطای توازن، خطای توقف و شروع کاذب در خطوط ۲۷۹ تا ۳۰۸ همان پرونده به‌صورت
خودبررسی نوشته شده‌اند:

\begin{dsdlcodeblock}
\begin{verbatim}
279  $display("TEST: required parity error");
280  base_count = raw_completion_count;
281  drive_raw_frame(SECOND_DATA, ~(^SECOND_DATA), 1'b1);
282  wait_raw_count(base_count + 1);
283  if (raw_correct_history[base_count] !== 1'b0 ||
284      raw_data_history[base_count] !== FIRST_DATA)
285      $fatal(1, "FAIL: parity error semantics");
286
287  $display("TEST: required stop-bit error");
288  base_count = raw_completion_count;
289  drive_raw_frame(SECOND_DATA, ^SECOND_DATA, 1'b0);
290  wait_raw_count(base_count + 1);
291  if (raw_correct_history[base_count] !== 1'b0 ||
292      raw_data_history[base_count] !== FIRST_DATA)
293      $fatal(1, "FAIL: stop-bit error semantics");
\end{verbatim}
\end{dsdlcodeblock}

آزمون شروع کاذب و بازیابی پس از آن در قطعهٔ بعدی آمده است:

\begin{dsdlcodeblock}
\begin{verbatim}
294
295  $display("TEST: false-start rejection and recovery");
296  base_count = raw_completion_count;
297  @(negedge clk);
298  raw_rx = 1'b0;
299  @(negedge clk);
300  raw_rx = 1'b1;
301  repeat (4 * BIT_TICKS) @(negedge clk);
302  if (raw_completion_count != base_count)
303      $fatal(1, "FAIL: false start completed a frame");
304  drive_raw_frame(SECOND_DATA, ^SECOND_DATA, 1'b1);
305  wait_raw_count(base_count + 1);
306  if (raw_correct_history[base_count] !== 1'b1 ||
307      raw_data_history[base_count] !== SECOND_DATA)
308      $fatal(1, "FAIL: recovery after false start");
\end{verbatim}
\end{dsdlcodeblock}

در آزمون مستقل همهٔ مقادیر، حلقهٔ خطوط ۴۵ تا ۶۶ هر کد هفت‌بیتی را ارسال می‌کند، پایان
دریافت را با مهلت محدود انتظار می‌کشد، داده و اعتبار را می‌سنجد و یک‌چرخه‌ای‌بودن دو پرچم
را نیز کنترل می‌کند:

\begin{dsdlcodeblock}
\begin{verbatim}
45  for (value = 0; value < 128; value = value + 1) begin
46      @(negedge clk);
47      send_data = value[6:0];
48      new_data = 1'b1;
49      @(negedge clk);
50      new_data = 1'b0;
51
52      cycles = 0;
53      while (rec_new_data !== 1'b1 && cycles < TIMEOUT_CYCLES) begin
54          @(negedge clk);
55          cycles = cycles + 1;
56      end
57      if (rec_new_data !== 1'b1)
58          $fatal(1, "FAIL: value %0d timed out", value);
59      if (correct_data !== 1'b1 || rec_data !== value[6:0])
60          $fatal(1, "FAIL: sent=%0d received=%0d correct=%b",
61                 value, rec_data, correct_data);
62
63      @(negedge clk);
64      if (rec_new_data !== 1'b0 || correct_data !== 1'b0)
65          $fatal(1, "FAIL: output pulse width for value %0d", value);
66  end
\end{verbatim}
\end{dsdlcodeblock}


\section{تحلیل شکل‌موج‌های شبیه‌سازی}

شکل‌های این بخش خروجی اجرای میزآزمون‌ها با \لر{Icarus Verilog} و مشاهده در
\لر{GTKWave} هستند. هر تحلیل، محرک دقیق، سیگنال‌های قابل مشاهده و نتیجهٔ کنترل خودکار را
در کنار هم بیان می‌کند.


\subsection{ارسال و دریافت معتبر}

رفتار فرستنده برای دادهٔ نمونه در شکل~\رجوع{شکل:فرستنده معتبر} دیده می‌شود.

\شروع{شکل}[H]
\تنظیم‌ازوسط
\centerimg{uart_valid_sender.png}{17cm}
\شرح{شکل‌موج فرستنده برای ارسال معتبر دادهٔ \کد{7'b1011001}}
\برچسب{شکل:فرستنده معتبر}
\پایان{شکل}

در شکل~\رجوع{شکل:فرستنده معتبر}، \کد{TEST\_DATA} برابر \کد{7'b1011001} یا
\کد{0x59} است. \کد{new\_data} برای یک چرخه فعال می‌شود؛ در همان پذیرش،
\کد{data\_latch} مقدار \کد{0x59} و \کد{parity\_latch} مقدار ۰ را نگه می‌دارند و
\کد{busy} فعال می‌شود. \کد{tx} ابتدا \کد{Start=0}، سپس \کد{Parity=0}، بعد دنبالهٔ
\کد{D0} تا \کد{D6} برابر ۱، ۰، ۰، ۱، ۱، ۰، ۱ و سرانجام \کد{Stop=1} را تولید
می‌کند. هر بخش با \کد{BIT\_TICKS=10} دقیقاً ده چرخهٔ ساعت پایدار است. تغییرات
\کد{bit\_index} مرز بیت‌های داده را مشخص می‌کنند و بازگشت \کد{state} به \لر{IDLE}
هم‌زمان با غیرفعال‌شدن \کد{busy} است. در نتیجه، قفل داده، محاسبهٔ توازن، ترتیب
\لر{LSB-first} و زمان‌بندی فرستنده صحیح‌اند.

بازسازی همان قاب در گیرنده در شکل~\رجوع{شکل:گیرنده معتبر} نمایش داده شده است.

\شروع{شکل}[H]
\تنظیم‌ازوسط
\centerimg{uart_valid_receiver.png}{17cm}
\شرح{شکل‌موج گیرنده برای دریافت معتبر دادهٔ \کد{7'b1011001}}
\برچسب{شکل:گیرنده معتبر}
\پایان{شکل}

در شکل~\رجوع{شکل:گیرنده معتبر}، \کد{rx} سیگنال واقعی خط و \کد{rx\_meta} و
\کد{rx\_sync} به‌ترتیب نسخه‌های یک و دو مرحله همگام‌شدهٔ آن هستند؛ بنابراین،
هم‌ترازنبودن دقیق سه موج عمدی و ناشی از دو فلیپ‌فلاپ همگام‌ساز است، نه تصادفی‌بودن نرخ
بیت. ماشین پس از مشاهدهٔ صفر وارد \لر{START} می‌شود، در میانهٔ بیت آن را تأیید می‌کند،
\کد{parity\_latch=0} را می‌گیرد و سپس \کد{data\_latch} را از \کد{D0} تا \کد{D6}
تکمیل می‌کند. در پایان، \کد{Stop=1} است و توازن محاسبه‌شدهٔ صفر با بیت ذخیره‌شده برابر
است؛ ازاین‌رو \کد{rec\_data} به \کد{0x59} تغییر می‌کند و \کد{rec\_new\_data} و
\کد{correct\_data} هم‌زمان برای یک چرخه بالا می‌روند. این موج بازسازی صحیح و
اعتبارسنجی نهایی قاب را اثبات می‌کند.


\subsection{خطای توازن}

واکنش گیرنده به بیت توازن خراب در شکل~\رجوع{شکل:خطای توازن} دیده می‌شود.

\شروع{شکل}[H]
\تنظیم‌ازوسط
\centerimg{uart_parity_error.png}{17cm}
\شرح{شکل‌موج دریافت \کد{0x2D} با بیت توازن معکوس}
\برچسب{شکل:خطای توازن}
\پایان{شکل}

در بخش نخست شکل~\رجوع{شکل:خطای توازن}، قاب \کد{0x59} معتبر دریافت می‌شود تا
\کد{rec\_data} مقدار شناخته‌شدهٔ \کد{0x59} داشته باشد. قاب دوم دادهٔ
\کد{7'b0101101} یا \کد{0x2D} را حمل می‌کند. این داده نیز چهار بیت یک دارد و توازن
صحیح آن صفر است، اما محرک خام عمداً بیت توازن دریافتی را ۱ می‌کند. پس از دریافت هفت بیت
و توقف صحیح، گیرنده پایان قاب را با پالس یک‌چرخه‌ای \کد{rec\_new\_data} اعلام می‌کند؛
ولی چون \کد{parity\_latch=1} با \لر{XOR} داده، یعنی ۰، برابر نیست،
\کد{correct\_data} فعال نمی‌شود و \کد{rec\_data} نیز \کد{0x59} باقی می‌ماند. پس
پالس پایان با پالس اعتبار یکسان نیست و خطای توازن از ورود دادهٔ خراب به خروجی جلوگیری
می‌کند.


\subsection{خطای بیت توقف}

اثر صفرکردن عمدی بیت توقف در شکل~\رجوع{شکل:خطای توقف} نشان داده شده است.

\شروع{شکل}[H]
\تنظیم‌ازوسط
\centerimg{uart_stop_error.png}{17cm}
\شرح{شکل‌موج دریافت \کد{0x2D} با بیت توقف صفر}
\برچسب{شکل:خطای توقف}
\پایان{شکل}

در شکل~\رجوع{شکل:خطای توقف} نیز قاب معتبر \کد{0x59} ابتدا خروجی را مقداردهی می‌کند. قاب
خام بعدی حاوی \کد{0x2D} و توازن صحیح صفر است، اما \کد{Stop} عمداً صفر نگه داشته
می‌شود. چون \کد{D6} دادهٔ \کد{0x2D} نیز صفر است، \کد{D6} و بیت توقف خراب در موج به
شکل یک بازهٔ صفر طولانی و بدون لبهٔ میان آن‌ها دیده می‌شوند. در پایان بازه،
\کد{rec\_new\_data} یک چرخه فعال می‌شود، اما شرط \کد{rx\_sync=1} برقرار نیست؛ در
نتیجه \کد{correct\_data} صفر و \کد{rec\_data} برابر \کد{0x59} باقی می‌ماند. حتی با
داده و توازن صحیح، \کد{Stop=0} قاب را نامعتبر می‌کند.


\subsection{قاب‌های متوالی}

ترتیب دریافت دو قاب الزامی در شکل~\رجوع{شکل:قاب‌های متوالی} قابل مشاهده است.

\شروع{شکل}[H]
\تنظیم‌ازوسط
\centerimg{uart_back_to_back.png}{17cm}
\شرح{شکل‌موج ارسال متوالی \کد{0x59} و \کد{0x2D}}
\برچسب{شکل:قاب‌های متوالی}
\پایان{شکل}

در شکل~\رجوع{شکل:قاب‌های متوالی} ابتدا \کد{0x59} و سپس \کد{0x2D} ارسال می‌شود. درخواست
دوم در نخستین فرصت قانونی پس از پایین‌آمدن \کد{busy} اعمال شده است. میان دو قاب یک بازهٔ
کوتاه \لر{IDLE} و \کد{busy=0} دیده می‌شود که اجازهٔ پذیرش دادهٔ دوم را می‌دهد؛ سپس
\کد{busy} دوباره برای قاب بعدی فعال می‌شود. گیرنده دو پالس \کد{rec\_new\_data} و
\کد{correct\_data} تولید می‌کند. نمونهٔ ثبت‌شدهٔ نخست \کد{0x59} و نمونهٔ دوم
\کد{0x2D} است. ترتیب دو مقدار و وجود دقیقاً دو تکمیل معتبر نشان می‌دهد قاب‌ها بدون
جابه‌جایی یا از‌دست‌رفتن داده پردازش می‌شوند.


\subsection{نادیده‌گرفتن درخواست هنگام اشتغال}

آزمون محافظت از قاب فعال در شکل~\رجوع{شکل:نادیده‌گرفتن busy} آمده است.

\شروع{شکل}[H]
\تنظیم‌ازوسط
\centerimg{uart_busy_ignore.png}{17cm}
\شرح{شکل‌موج درخواست \کد{0x2D} هنگام ارسال فعال \کد{0x59}}
\برچسب{شکل:نادیده‌گرفتن busy}
\پایان{شکل}

در شکل~\رجوع{شکل:نادیده‌گرفتن busy} ارسال \کد{0x59} آغاز می‌شود و
\کد{data\_latch} همین مقدار را نگه می‌دارد. در میانهٔ \کد{busy}،
\کد{send\_data} به \کد{0x2D} تغییر می‌کند و \کد{new\_data} دوباره برای یک چرخه فعال
می‌شود. چون منطق پذیرش فقط در \لر{IDLE} قرار دارد، ثبات داده تغییر نمی‌کند و \کد{tx}
ادامهٔ همان قاب \کد{0x59} را نمایش می‌دهد. در پایان تنها یک پالس دریافت معتبر با
\کد{rec\_data=0x59} دیده می‌شود و پس از زمان کافی هیچ پالس دومی ایجاد نمی‌گردد. پس
درخواست هنگام اشتغال نه قاب جاری را مخدوش می‌کند و نه یک قاب ناخواسته در صف می‌سازد.


\subsection{بازنشانی و حالت بیکار}

مقادیر تعریف‌شدهٔ مدار در بازنشانی و بیکاری در شکل~\رجوع{شکل:بازنشانی و بیکاری} دیده
می‌شوند.

\شروع{شکل}[H]
\تنظیم‌ازوسط
\centerimg{uart_reset_idle.png}{17cm}
\شرح{شکل‌موج بازنشانی فعال-پایین و پایداری حالت بیکار}
\برچسب{شکل:بازنشانی و بیکاری}
\پایان{شکل}

در شکل~\رجوع{شکل:بازنشانی و بیکاری}، \کد{send\_data} از ابتدا \کد{0x7F} است، اما
\کد{new\_data} صفر نگه داشته می‌شود. \کد{rstN} بازنشانی فعال-پایین را تقریباً از ۱ تا
۲۰ نانوثانیه اعمال می‌کند. مقادیر قرمز \لر{X} در لحظهٔ نخست، وضعیت طبیعی پیش از
اثرکردن بازنشانی هستند. با فعال‌شدن بازنشانی، حالت، شمارنده‌ها و ثبات‌ها صفر می‌شوند،
\کد{rx\_meta} و \کد{rx\_sync} به ۱ می‌رسند، \کد{tx=1}، \کد{busy=0}،
\کد{rec\_data=0} و خروجی‌های پالسی صفر می‌شوند. پس از آزادشدن \کد{rstN} تا پایان بازهٔ
۴۰ نانوثانیه، مدار در \لر{IDLE} پایدار می‌ماند. ثابت‌بودن \کد{send\_data} به‌تنهایی
انتقالی آغاز نمی‌کند و نیاز به \کد{new\_data} تأیید می‌شود.


\subsection{پالس‌های خروجی گیرنده}

پهنای اعلان پایان و اعتبار در شکل~\رجوع{شکل:پالس‌های خروجی} مقایسه شده است.

\شروع{شکل}[H]
\تنظیم‌ازوسط
\centerimg{uart_output_pulses.png}{17cm}
\شرح{مقایسهٔ پالس‌های خروجی برای قاب معتبر و قاب دارای خطای توازن}
\برچسب{شکل:پالس‌های خروجی}
\پایان{شکل}

در شکل~\رجوع{شکل:پالس‌های خروجی} نخست یک انتقال بازگشتی معتبر \کد{0x59} انجام می‌شود.
نزدیک ۱ میکروثانیه، \کد{rec\_data=0x59} و \کد{rec\_new\_data} و
\کد{correct\_data} دقیقاً برای یک دورهٔ \کد{clk} فعال می‌شوند و در لبهٔ بعدی پایین
می‌آیند. سپس ورودی خام همان دادهٔ \کد{0x59} را با بیت توازن معکوس ۱ اعمال می‌کند. نزدیک
۲٫۰۳ میکروثانیه، \کد{raw\_rec\_new\_data} برای یک دورهٔ ساعت فعال می‌شود، اما
\کد{raw\_correct\_data} همواره صفر و \کد{raw\_rec\_data} برابر صفر باقی می‌ماند.
این مقایسه نشان می‌دهد اعلان پایان برای هر قاب کامل ایجاد می‌شود، در حالی که اعلان اعتبار
فقط به قاب سالم تعلق دارد و هر دو خروجی ماهیت پالسی یک‌چرخه‌ای دارند.


\subsection{شروع کاذب و بازیابی}

رد پالس کوتاه خط و دریافت صحیح بعدی در شکل~\رجوع{شکل:شروع کاذب} نمایش داده شده است.

\شروع{شکل}[H]
\تنظیم‌ازوسط
\centerimg{uart_false_start.png}{17cm}
\شرح{شکل‌موج رد شروع کاذب ۱۰ نانوثانیه‌ای و بازیابی با \کد{0x59}}
\برچسب{شکل:شروع کاذب}
\پایان{شکل}

با \کد{BIT\_TICKS=10} و دورهٔ ساعت ۱۰ نانوثانیه، نیمهٔ زمان بیت ۵۰ نانوثانیه است. در
شکل~\رجوع{شکل:شروع کاذب}، \کد{rx} در حدود ۵۰ نانوثانیه فقط برای یک چرخه، یعنی ۱۰
نانوثانیه، صفر می‌شود. پس از عبور پالس از \کد{rx\_meta} و \کد{rx\_sync}، ماشین ممکن
است موقتاً وارد \لر{START} شود، اما هنگام تأیید میانی خط دوباره ۱ است و ماشین مستقیماً به
\لر{IDLE} بازمی‌گردد. تا این مرحله نه اعلان پایان و نه اعتبار فعال نمی‌شوند و شمار تکمیل‌ها
صفر می‌ماند. قاب معتبر بعدی با دادهٔ \کد{0x59} به‌طور کامل دریافت می‌شود،
\کد{rec\_data=0x59} و هر دو پرچم برای یک چرخه فعال می‌شوند و
\کد{completion\_count} به ۱ می‌رسد. پس گیرنده هم شروع کوتاه کاذب را رد می‌کند و هم بدون
بازنشانی بازیابی می‌شود.


\subsection{پیمایش تمام مقادیر هفت‌بیتی}

آغاز پیمایش در شکل~\رجوع{شکل:همه مقادیر ۱} دیده می‌شود.

\شروع{شکل}[H]
\تنظیم‌ازوسط
\centerimg{uart_all_values_1.png}{17cm}
\شرح{بخش نخست شکل‌موج پیمایش تمام مقادیر هفت‌بیتی}
\برچسب{شکل:همه مقادیر ۱}
\پایان{شکل}

شکل~\رجوع{شکل:همه مقادیر ۱} آغاز پیمایش را نشان می‌دهد. در پنجرهٔ قابل مشاهده،
\کد{send\_data} تقریباً از \کد{0x00} تا \کد{0x19} پیش می‌رود و \کد{rec\_data} با
یک قاب تأخیر، مقادیر تقریباً \کد{0x00} تا \کد{0x18} را دنبال می‌کند. برای هر مقدار، یک
پالس \کد{rec\_new\_data} و یک پالس \کد{correct\_data} دیده می‌شود.

ادامهٔ نخست دنباله در شکل~\رجوع{شکل:همه مقادیر ۲} آمده است.

\شروع{شکل}[H]
\تنظیم‌ازوسط
\centerimg{uart_all_values_2.png}{17cm}
\شرح{بخش دوم شکل‌موج پیمایش تمام مقادیر هفت‌بیتی}
\برچسب{شکل:همه مقادیر ۲}
\پایان{شکل}

شکل~\رجوع{شکل:همه مقادیر ۲} ادامهٔ پیمایش را در حدود
\کد{send\_data=0x18..0x31} و \کد{rec\_data=0x17..0x30} نمایش می‌دهد. افزایش مرتب
داده و جفت‌شدن هر دریافت با پالس اعتبار نشان می‌دهد مرز قاب‌ها در دنبالهٔ طولانی حفظ شده
است.

بخش میانی دنباله در شکل~\رجوع{شکل:همه مقادیر ۳} دیده می‌شود.

\شروع{شکل}[H]
\تنظیم‌ازوسط
\centerimg{uart_all_values_3.png}{17cm}
\شرح{بخش سوم شکل‌موج پیمایش تمام مقادیر هفت‌بیتی}
\برچسب{شکل:همه مقادیر ۳}
\پایان{شکل}

شکل~\رجوع{شکل:همه مقادیر ۳} بخش میانی آزمون را در حدود
\کد{send\_data=0x2D..0x46} و \کد{rec\_data=0x2C..0x45} نشان می‌دهد. مقادیر دارای
الگوهای متفاوت صفر و یک بدون خطای توازن یا ترتیب بیت بازسازی می‌شوند.

آخرین پنجرهٔ موجود در شکل~\رجوع{شکل:همه مقادیر ۴} آمده است.

\شروع{شکل}[H]
\تنظیم‌ازوسط
\centerimg{uart_all_values_4.png}{17cm}
\شرح{بخش چهارم شکل‌موج پیمایش تمام مقادیر هفت‌بیتی}
\برچسب{شکل:همه مقادیر ۴}
\پایان{شکل}

شکل~\رجوع{شکل:همه مقادیر ۴} ادامهٔ آزمون را در حدود
\کد{send\_data=0x46..0x60} و \کد{rec\_data=0x46..0x5F} نشان می‌دهد. این چهار تصویر
فقط پنجره‌های پیوسته‌ای از اجرای طولانی هستند؛ پوشش کامل از خود حلقهٔ میزآزمون و نتیجهٔ
خودبررسی آن حاصل می‌شود. حلقه مقدار \کد{value} را از ۰ تا ۱۲۷ تغییر می‌دهد، برای هر مقدار
انتقال را آغاز می‌کند، منتظر \کد{rec\_new\_data} می‌ماند و برابری \کد{rec\_data} با
مقدار ارسالی و فعال‌بودن \کد{correct\_data} را کنترل می‌کند. بنابراین، همهٔ ۱۲۸ کد
هفت‌بیتی از \کد{0x00} تا \کد{0x7F} آزمایش شده‌اند، هرچند تمام آن‌ها هم‌زمان در یک تصویر
جا نمی‌گیرند.


\section{خروجی متنی شبیه‌ساز}

میزآزمون تجمیعی با دستورهای \کد{\$display} نام هر مرحله و نتیجهٔ نهایی را چاپ می‌کند.
خروجی کامل اجرای موفق به‌صورت زیر است:

\begin{dsdlcodeblock}
\begin{verbatim}
PASS: reset and idle outputs
TEST: exact valid frame and symbol timing
TEST: request asserted while busy
TEST: required consecutive frames
TEST: exhaustive seven-bit loopback
TEST: valid raw receiver frame
TEST: required parity error
TEST: required stop-bit error
TEST: false-start rejection and recovery
PASS: complete DSDL7 UART regression
\end{verbatim}
\end{dsdlcodeblock}

این خروجی، \لر{Transcript} متنی مورد نیاز را فراهم می‌کند. رسیدن اجرا به خط پایانی
\لر{PASS} به این معناست که هیچ‌یک از کنترل‌های خطای داخلی میزآزمون فعال نشده‌اند. نتایج
موج و خروجی متنی یکدیگر را تکمیل می‌کنند: شکل‌موج ترتیب زمانی را نشان می‌دهد و کنترل
خودکار، مقدار و پهنای پالس را در لبه‌های ساعت بررسی می‌کند.


\section{سنتز و بررسی ساختار در \لر{Quartus}}

خلاصهٔ نتیجه و منابع سنتز در
شکل~\رجوع{شکل:خلاصه سنتز Quartus} آمده است.

\شروع{شکل}[H]
\تنظیم‌ازوسط
\centerimg{quartus3.jpg}{17cm}
\شرح{خلاصهٔ موفق \لر{Analysis \& Synthesis} و مصرف منابع پیمانهٔ \کد{UARTTop}}
\برچسب{شکل:خلاصه سنتز Quartus}
\پایان{شکل}

شکل~\رجوع{شکل:خلاصه سنتز Quartus} صفحهٔ \لر{Analysis \& Synthesis Summary} را پس از
اجرای موفق سنتز نشان می‌دهد. وضعیت \لر{Analysis \& Synthesis} برابر
\لر{Successful}، نام بازبینی \کد{q7}، پیمانهٔ سطح بالا \کد{UARTTop} و خانوادهٔ قطعه
\لر{Cyclone IV GX} است. برای طرح ۱۰۶ عنصر منطقی\LTRfootnote{Logic Elements}، ۹۱ تابع
ترکیبی\LTRfootnote{Combinational Functions}، ۶۱ ثبات
اختصاصی، ۶۱ ثبات در مجموع و ۲۱
پایه گزارش شده است. تعداد بیت‌های حافظه، ضرب‌کننده‌های ۹ بیتی،
\لر{PLL}ها و منابع فرستنده‌گیرندهٔ پرسرعت همگی صفر هستند. بنابراین، این
تصویر هم موفقیت سنتز و هم هزینهٔ واقعی منابع طرح را مستند می‌کند.

ساختار سلسله‌مراتبی سطح بالا در
شکل~\رجوع{شکل:نمای سلسله‌مراتبی RTL} دیده می‌شود.

\شروع{شکل}[H]
\تنظیم‌ازوسط
\centerimg{quartus4.jpg}{17cm}
\شرح{نمای سلسله‌مراتبی \لر{RTL} از فرستنده، گیرنده و مسیر بازگشتی داخلی}
\برچسب{شکل:نمای سلسله‌مراتبی RTL}
\پایان{شکل}

شکل~\رجوع{شکل:نمای سلسله‌مراتبی RTL} دو بلوک \کد{UARTSender} و \کد{UARTReceiver} را
به‌وضوح جدا نشان می‌دهد. ورودی‌های \کد{clk} و \کد{rstN} به هر دو بلوک می‌رسند و
\کد{new\_data} و \کد{send\_data[6:0]} ورودی‌های فرستنده‌اند. خروجی \کد{tx} فرستنده
هم به درگاه خروجی \کد{tx} و هم مستقیماً به ورودی \کد{rx}
گیرنده متصل شده است.
خروجی‌های \کد{busy}، \کد{rec\_data[6:0]}، \کد{rec\_new\_data} و
\کد{correct\_data} نیز از بلوک مربوط به خود به درگاه‌های سطح بالا می‌رسند. بنابراین، نما
اتصال بازگشتی داخلی مورد نیاز آزمایش را تأیید می‌کند.

نگاشت تفصیلی شبکهٔ سنتزشده در
شکل~\رجوع{شکل:شبکه ورودی خروجی Quartus} ارائه شده است.

\شروع{شکل}[H]
\تنظیم‌ازوسط
\centerimg{quartus2.jpg}{17cm}
\شرح{نمای تفصیلی شبکهٔ سنتزشده و بافرهای ورودی و خروجی پیمانهٔ \کد{UARTTop}}
\برچسب{شکل:شبکه ورودی خروجی Quartus}
\پایان{شکل}

شکل~\رجوع{شکل:شبکه ورودی خروجی Quartus} ده ورودی شامل \کد{clk}، هفت بیت
\کد{send\_data}، \کد{rstN} و \کد{new\_data} را نشان می‌دهد که با بافرهای
ورودی از نوع \کد{IO\_IBUF} وارد ساختار می‌شوند. در سمت
خروجی نیز هفت بیت \کد{rec\_data} به‌همراه
\کد{correct\_data}، \کد{rec\_new\_data}، \کد{busy} و \کد{tx} از یازده بافر
خروجی از نوع \کد{IO\_OBUF} عبور می‌کنند؛ در نتیجه مجموع
۲۱ پایه با خلاصهٔ سنتز یکسان است. تفکیک
بیت‌های دو گذرگاه هفت‌بیتی و حفظ دو بلوک فرستنده و گیرنده، نگاشت درست
رابط سطح بالا را
نشان می‌دهد.

نحوهٔ تشخیص ماشین حالت گیرنده در شکل~\رجوع{شکل:کدگذاری حالت گیرنده Quartus} ثبت شده است.

\شروع{شکل}[H]
\تنظیم‌ازوسط
\centerimg{quartus1.jpg}{17cm}
\شرح{تشخیص پنج حالت گیرنده و کدگذاری یک‌داغ آن‌ها در گزارش \لر{State Machines}}
\برچسب{شکل:کدگذاری حالت گیرنده Quartus}
\پایان{شکل}

شکل~\رجوع{شکل:کدگذاری حالت گیرنده Quartus} گزارش \لر{State Machines} را برای ثبات
\کد{state} نمونهٔ \کد{UARTReceiver} نشان می‌دهد. \لر{Quartus} پنج حالت
\لر{IDLE}، \لر{START}، \لر{PARITY}، \لر{DATA} و \لر{STOP} را تشخیص داده و
\لر{Encoding Type} را \لر{One-Hot} انتخاب کرده است؛ در هر ردیف تنها بیت متناظر با همان
حالت مقدار ۱ دارد. هرچند حالت‌ها در کد \لر{Verilog} با مقادیر سه‌بیتی تعریف شده‌اند، ابزار
سنتز آن‌ها را برای پیاده‌سازی به کدگذاری یک‌داغ\LTRfootnote{One-Hot Encoding} تبدیل کرده
است، بی‌آنکه ترتیب یا رفتار ماشین
حالت تغییر کند. در درخت سمت چپ، ماشین حالت فرستنده نیز شناسایی شده است، اما جدول بازشده
مشخصاً به گیرنده تعلق دارد.


\section{تحلیل نهایی نتایج}

برای دادهٔ \کد{0x59}، خروجی سریال دقیقاً مطابق ترتیب \لر{Start}، \لر{Parity}،
\کد{D0} تا \کد{D6} و \لر{Stop} ساخته شد و گیرنده همان \کد{0x59} را همراه پالس
اعتبار بازسازی کرد. آزمون قالب با \کد{BIT\_TICKS=434} زمان واقعی هر بیت را کنترل می‌کند،
در حالی که شکل‌موج‌های \کد{BIT\_TICKS=10} همان منطق را با وضوح دیداری بیشتر نشان
می‌دهند. بنابراین، کاهش پارامتر در شبیه‌سازی تغییر پروتکل نیست، بلکه تنها فشرده‌سازی مقیاس
زمانی است.

آزمون‌های توازن و توقف نشان دادند شرط اعتبار از دو کنترل مستقل تشکیل شده است. در هر دو
خطا، تکمیل قاب اعلام شد، اما \کد{correct\_data} صفر ماند و \کد{rec\_data} به مقدار
نامعتبر تغییر نکرد. این رفتار امکان تمایز میان «قاب پردازش‌شده» و «دادهٔ قابل استفاده» را
فراهم می‌کند. آزمون \کد{busy} نیز ثابت کرد قفل ورودی، قاب فعال را در برابر تغییرات
\کد{send\_data} و درخواست اضافی محافظت می‌کند. آزمون شروع کاذب نقش نمونه‌برداری میانی را
نشان داد و پیمایش ۱۲۸ مقدار، عملکرد را فراتر از یک بردار نمونه پوشش داد.


\section{نتیجه‌گیری}

فرستنده و گیرندهٔ \لر{UART} هفت‌بیتی مطابق قالب ویژهٔ آزمایش طراحی شده‌اند. ماشین حالت
فرستنده قاب را با توازن \لر{XOR} و ترتیب \لر{LSB-first} تولید می‌کند و \کد{busy}
محدودهٔ اشتغال را دقیقاً مشخص می‌سازد. گیرنده ورودی ناهمزمان را همگام، شروع را در میانهٔ
بیت تأیید، داده را بازسازی و توازن و توقف را کنترل می‌کند. اتصال بازگشتی سطح بالا و مجموعه
آزمون‌های خودکار، دریافت صحیح، رد قاب خراب، حفظ دادهٔ معتبر، رفتار بازنشانی، پالس‌های
یک‌چرخه‌ای، قاب‌های متوالی، رد درخواست هنگام اشتغال، بازیابی پس از شروع کاذب و تمام فضای
دادهٔ هفت‌بیتی را پوشش می‌دهند. در نتیجه، رفتار منطقی و زمانی طراحی با الزامات آزمایش هفتم
سازگار است.
